\documentclass[compress]{beamer}
\input{../../resources/preamble.sty}
\addbibresource{../../resources/literature.bib}
\graphicspath{{../../resources/img/}}


\begin{document}

\title[Big Data and Automated Content Analysis]{\textbf{Big Data and Automated Content Analysis (6EC)} 
\\Week 3: »Data Wrangling and Exploratory Analysis«
\\Monday}
\author[Anne Kroon]{Anne Kroon\\ \footnotesize{a.c.kroon@uva.nl \\}}
\date{November 10, 2025}
\institute[UvA CW]{UvA RM Communication Science}


\begin{frame}{}
	\titlepage
\end{frame}

\begin{frame}{Today}
	\tableofcontents
\end{frame}

\question{Everything clear from last weeks?}

\begin{frame}{Main points from last week}

\begin{alertblock}{I assume that by now, everybody knows:}
\begin{itemize}
\item how to read and write different types of files
\item how to work with both nested and tabular data
\item how to deal with data of an unknown structure (such as responses from an unfamiliar JSON API)
\item \ldots and of course all the basics: data types, for loops, conditions, etc.
\end{itemize}
\end{alertblock}
\end{frame}

\begin{frame}[standout]
This week, we will look into tabular, relatively traditional data. But remember what we said last week about different data structures -- we will work with them a lot as well in the upcoming weeks!
\end{frame}

\input{../../modules/basics/pandas.tex}

\input{../../modules/basics/visualization.tex}


\question{Any questions?}

\section{Next steps}

\begin{frame}[standout]

Make sure you understood today’s concepts.

Re-read the chapters and play with the API example:  
\url{https://github.com/uvacw/teaching-bdaca/blob/main/6ec-course/week02/API-getting-started.ipynb}

Try changing the query, looping over results, and inspecting data types.

Exercises for this week:  
\url{https://github.com/uvacw/teaching-bdaca/tree/main/6ec-course/week03/exercises-03-2}

We’ll work on them together on Thursday.
\end{frame}

\begin{frame}[allowframebreaks,plain]
  \printbibliography
\end{frame}




\end{document}
